\chapter{Introduction}\label{ch01:introduction}

\section{Motivation and setting}\label{ch01:motivation}

Every research community that models complex adaptive behavior eventually
confronts the same methodological hazard: the informal diagram. A block sketch
of a hypothesized cognitive pathway --- boxes for stages, arrows for influence,
annotations for thresholds and gains --- is cheap to draw, easy to circulate,
and almost impossible to audit. Two readers of the same sketch can disagree
about whether an arrow denotes causation, correlation, data flow, or mere
suggestion; whether a threshold is a design intention or a measured property;
whether a feedback loop is stable, unstable, or simply unanalyzed; and whether
the sketch as a whole is a claim about the world or a claim about a model of
the world. The sketch has no type system, no versioning, no diagnostics, no
replayable execution semantics, and no evidence discipline. When such sketches
accumulate into a large body of design work --- in the present case, a
743-page engineering baseline --- the hazard compounds: assertions of very
different epistemic standing sit side by side on the page, typeset with equal
confidence, and nothing in the medium itself distinguishes a tested software
behavior from a hypothesis, an idealization, or an error.

Hardware engineering solved a structurally similar problem decades ago, not by
proving anything about physics that physicists did not already know, but by
changing the medium of design description. Machine-readable hardware
description languages, explicit state formalisms, netlist-centered design
databases, design-rule checkers that fail closed, deterministic simulators
with archived traces, and sign-off disciplines that separate verification from
validation together converted circuit design from a craft of annotated
drawings into an auditable engineering practice. Chapter~\ref{ch02:background}
surveys that heritage in detail. The essential observation for this monograph
is that the value of those tools was largely \emph{methodological}: they made
design intent explicit, mechanically checkable, and reproducible, and they
made whole categories of error observable before fabrication.

CAS PCBA 8 --- the subject of this monograph --- asks whether the same
methodological move is available to cognitive-system modeling. It proposes a
typed, provenance-bearing research architecture in which a candidate
cognitive-system hypothesis is written down as a \emph{circuit-shaped model}:
a directed, typed network of components with declared ports, signal kinds,
units, parameters, state annotations, calibration references, feedback
declarations, and evidence limits, subject to a deterministic design-rule
checker and, where supported, a narrow deterministic simulator that emits
canonical traces. The proposal is emphatically \emph{not} that cognition is
electronics. The proposal is that the engineering \emph{organization} of
electronics --- its description discipline, its checking discipline, and its
evidence discipline --- can be borrowed as a modeling grammar while every
physical, causal, safety, and product claim is explicitly withheld.

This introductory chapter states the research question precisely, explains
what ``circuit'' does and does not mean in this work, develops the three
load-bearing separations that keep the analogy from decaying into an identity
claim, enumerates the six contributions, states the qualified-affirmative
thesis defended in the remainder of the monograph, fixes the claim discipline
under which every subsequent chapter operates, and closes with a reading
guide.

\section{The name of the discipline}\label{ch01:name}

The monograph's subtitle names its subject \emph{adaptive cognitive circuit
engineering}, and each word of the name is a scoped commitment worth fixing
before the technical development begins.

\emph{Engineering} commits the work to the practices of that word: declared
interfaces, versioned artifacts, deterministic tools, catalogs of checks,
archived evidence, and an explicit boundary between what has been tested and
what has been asserted. It deliberately does not commit the work to science:
no chapter of this monograph reports an experiment, and the word
``engineering'' is chosen partly to keep that fact in the reader's view.
What is engineered here is a modeling and checking apparatus --- a language,
a checker, a simulator, adapters, and an evidence discipline --- and the
apparatus is the deliverable.

\emph{Circuit} commits the work to the specific engineering organization of
Definition~\ref{ch01:def:circuit} and Boundary~\ref{ch01:bnd:circuit}: typed
components, ports, nets, state annotations, calibration references,
deterministic traces. The commitment is organizational, never physical.

\emph{Cognitive} commits the work to a subject matter --- hypotheses about
cognitive systems, drawn from the CAS framework whose semantic identities
the typed model carries --- while committing it to no position on whether
any such hypothesis is true. The cognitive content enters the architecture
only through typed semantic references with their own namespaces, sources,
and definition hashes, and the architecture's machinery is deliberately
indifferent to whether those references denote anything real: it checks
structure, connectivity, calibration shape, and reachability identically
whether the board models a mind, a control loop, or nothing at all. That
indifference is a feature. A checker that needed the cognitive hypotheses
to be true in order to run would have smuggled the conclusion into the
tool.

\emph{Adaptive} commits the work, prospectively, to systems whose
parameters, structure, or behavior change over time --- the eventual
subject matter includes feedback, state, drift, and recalibration --- while
the honest present tense is narrower: the current candidate supports
declared feedback with unit delays and conservative gain bounds, and
treats stateful, stochastic, and richer adaptive semantics as explicitly
open, failing closed wherever they are encountered. The word marks the
direction of the research program, not an achieved capability.

The workflow the name describes is a loop the later chapters will make
precise. A hypothesis is written as a board (design); the board is checked
against the closed diagnostic catalog (rule check); supported boards are
simulated under declared profiles, emitting canonical traces (simulate);
boards, reports, and traces are audited against their sources and against
the claim discipline (audit); and the resulting artifacts enter the
evidence ledgers typed as what they are --- software receipts, never more
(evidence). Every stage is deterministic, identity-bound, and fail-closed,
and the loop's product is not a validated model of cognition but a
falsifiable, replayable, honestly bounded description of one.

\section{The research question}\label{ch01:question}

The engineering question is not whether cognition \emph{is} electronics. It
is:

\begin{quote}
Can a circuit-shaped modeling language make cognitive-system hypotheses more
explicit, type-safe, replayable, and falsifiable without converting an analogy
into a physical, causal, safety, or product claim?
\end{quote}

Each adjective in the question carries weight, and it is worth unpacking them
before anything else is asserted.

\emph{Explicit} means that every element of a hypothesis that the modeler
relies on --- the components, their interfaces, the values they may carry, the
transfer behavior they are assumed to exhibit, their statefulness, their time
semantics, their calibration status, and the boundary of what the model does
not claim --- appears in the model text as a declared, typed object rather
than as an unstated assumption. An informal diagram is explicit about
topology and silent about almost everything else; a circuit-shaped model in
the sense of this monograph is required to be explicit about all of it, and a
model that omits a required declaration does not degrade gracefully into a
vaguer model --- it fails the checker.

\emph{Type-safe} means that the modeling language distinguishes kinds of
things that must not be silently interchanged, and that the distinction is
enforced mechanically. A semantic state literal is not a probability profile
over states; a probability profile is not an element of a finite kernel; a
design parameter is not a measurement; a locally asserted calibration label is
not authenticated calibration evidence; an integer $1$, a Boolean
\emph{true}, and a string \texttt{"true"} are three distinct values
everywhere. Type safety in this sense is the primary technical instrument by
which the analogy is prevented from leaking: the places where an analogy
would like to become an identity are exactly the places where the type system
refuses a coercion.

\emph{Replayable} means that every executable claim is bound to enough
identity material --- source hashes, profile identities, decimal-context
identity, component-catalog identity, input samples, and a canonical trace ---
that an independent party holding the same artifacts can rerun the exact
computation and compare results value for value. Replayability is a property
of the software contract, not of the world: replaying a simulation
demonstrates that the software is deterministic under its declared profile,
and nothing more.

\emph{Falsifiable} means that the model text exposes surfaces on which it can
fail, and that failures are preserved rather than repaired away. The
design-rule checker can reject a board; the threshold-reachability analysis
can prove that a published safety stage cannot activate; a grammar-coverage
audit can show that a published formal grammar does not cover its own
published corpus. The monograph's source audit (summarized in
Chapter~\ref{ch03:source} and developed in full in later chapters) records
precisely such failures against the engineering baseline itself, and the
architecture retains them as first-class findings and permanent regression
fixtures.

The final clause of the question --- \emph{without converting an analogy into
a physical, causal, safety, or product claim} --- is the constraint that
makes the question nontrivial. It is easy to build an evocative circuit
metaphor; the literature of every era is full of cognitive models dressed in
the technological costume of their day --- hydraulic, telegraphic,
telephonic, computational. The historical failure mode of such metaphors is
not that they are useless but that they are boundless: the costume is
mistaken for anatomy. The question this monograph addresses is whether the
costume can be worn with a mechanically enforced boundary --- whether the
analogy can be made load-bearing for design, audit, and falsification while
remaining, checkably and permanently, an analogy.

\section{What ``circuit'' means here --- and what it does not}\label{ch01:circuit}

\begin{definition}[Circuit-shaped model, informal]\label{ch01:def:circuit}
In this monograph, a \emph{circuit-shaped model} is a directed, typed network
whose elements must declare: (i) interfaces (ports with direction, signal
kind, and unit); (ii) allowed values (typed parameters and finite-state
annotations); (iii) transfer behavior (a declared component profile, where
one is supported); (iv) statefulness and time semantics; (v) calibration
references for any point at which observation is claimed; and (vi) evidence
limits (a fixed research boundary carried in the model's own identity
material). ``Circuit'' names this engineering organization and nothing else.
\end{definition}

\begin{boundary}[What ``circuit'' does not assert]\label{ch01:bnd:circuit}
The word ``circuit'' in CAS PCBA 8 does not assert a physical electrical
implementation. It does not identify any model quantity with voltage,
current, resistance, charge, flux, or power in the SI sense. It does not
assert that any cognitive system --- human, artificial, or biological ---
contains components corresponding to the model's components, nets
corresponding to its nets, or thresholds corresponding to its thresholds. It
does not assert that a simulated behavior of the model is a behavior of any
system outside the model. Any statement in this monograph that appears to
approach such an assertion is governed by this boundary and by the claim
quarantine of Section~\ref{ch01:discipline}.
\end{boundary}

The choice of the circuit organization, rather than some other engineering
organization, is motivated by a specific list of disciplines that circuit
practice developed and that informal cognitive diagrams lack. Six are worth
naming here, because each reappears as a concrete mechanism in later
chapters.

First, \emph{interface discipline}: in circuit practice a component is
usable only through declared ports, and connecting incompatible ports is a
detectable error rather than a silent reinterpretation. Second,
\emph{connectivity discipline}: a netlist makes the complete connection
structure of a design a single queryable object, so that floating inputs,
contending drivers, and dangling outputs are mechanical findings. Third,
\emph{rule discipline}: design-rule checking established, industrially, that
a design is not done when it looks right but when it passes a versioned,
deterministic body of checks --- and that the checks fail closed. Fourth,
\emph{trace discipline}: simulation practice archives waveforms against exact
design versions, so that a claimed behavior is inseparable from the identity
of the thing that behaved. Fifth, \emph{calibration discipline}: a probe
reading means nothing without an instrument identity, a protocol, and an
uncertainty; hardware measurement culture keeps design values and measured
values in separate ledgers. Sixth, \emph{sign-off discipline}: verification
(does the artifact implement its specification?) and validation (is the
specification adequate for an external purpose?) are distinct activities with
distinct evidence, and passing the first is never silently promoted into the
second.

Each of these is an organizational and methodological achievement, not a
physical one. None of them depends on electrons. That is precisely why they
are candidates for transfer: a modeling language for cognitive-system
hypotheses can adopt interface, connectivity, rule, trace, calibration, and
sign-off discipline wholesale, without adopting a single claim about
physics. The engineering literature surveyed in
Chapter~\ref{ch02:background} supports this transfer at the level of
\emph{method} --- machine-readable formal design description \cite{vhdl},
explicit state semantics \cite{harel}, verification as distinct from
validation \cite{sargent} --- while validating nothing about CAS itself.

\section{Three load-bearing separations}\label{ch01:separations}

The architecture's claim discipline rests on three separations. They are
called \emph{load-bearing} because every mechanism described in this
monograph --- the type system, the design-rule checker, the simulator's trace
bindings, the calibration rules, the adapter boundaries --- exists to enforce
at least one of them, and because the historical failure modes of
circuit-flavored cognitive modeling are, on inspection, violations of exactly
these three. Each separation is stated as a rule and then developed through
an extended examination of the failure it prevents.

\subsection{Separation I: model versus world}\label{ch01:sep:modelworld}

\begin{ruleenv}[Model versus world]\label{ch01:rule:modelworld}
A simulated transfer rule is behavior of the declared software profile, not
evidence that a cognitive or physical system obeys it. No output of the
architecture --- no green test, no deterministic trace, no diagnostic-free
design-rule report --- is admissible as evidence about any system other than
the software that produced it.
\end{ruleenv}

The failure this separation prevents is the quiet migration of a claim from
the model column to the world column. The migration rarely happens in a
single dishonest step; it happens through a chain of individually plausible
restatements. Consider a concrete, deliberately detailed example. A modeler
constructs a board containing a component intended to represent an
aversion-triggered refusal stage: an input carrying an activation-like
quantity, a comparator against a threshold parameter, an output that gates a
downstream path. The modeler writes a simulation fixture in which the input
ramps past the threshold and the gate closes, runs it, and obtains a green
result with a canonical trace. So far every statement is safe: \emph{the
software performed the declared typed transition and produced the declared
trace for the exact fixture under the exact profile}.

Now watch the restatements. In a status report, the green fixture becomes
``the refusal mechanism works.'' In a slide, ``the refusal mechanism works''
becomes ``the model refuses under overload.'' In a discussion section, ``the
model refuses under overload'' becomes ``this architecture exhibits the
refusal behavior observed in deployed systems,'' and in a grant abstract the
architecture has become evidence about cognition. Not one of the
restatements was flagged as a new claim; each merely dropped a
qualifier. The end state is a world-claim resting entirely on a software
fixture that was never designed, powered, or entitled to support it. The
history of simulation-based fields shows this slope is not hypothetical:
wherever simulation output is cheap and empirical access is expensive, the
simulated column bleeds into the empirical column unless the medium itself
resists.

A second example shows the same failure in the opposite emotional
register --- not enthusiasm but reassurance. Suppose a board's design-rule
report comes back clean: no errors, no warnings, every net driven, every
threshold reachable, every declared feedback loop bounded. The clean report
is a real and useful fact about the model text. The migration begins when
``the design checks out'' is heard as ``the design is safe,'' and completes
when a clean DRC report on a refusal-pathway model is offered as comfort
about the behavior of a deployed system that the model was never validated
against. Here the world-claim is not even about the model's subject; it is
about a different artifact entirely, connected to the model by nothing but
shared vocabulary. The checker's own fail-closed design anticipates the
first half of this slope --- unresolved analyses yield unresolved
diagnostics, never a pass --- but no tool can prevent a reader from
over-reading a clean report about A as evidence about B. Only the evidence
discipline can, by typing the report as a statement about the exact board
digest it names and nothing else.

CAS PCBA 8 makes the medium resist in three ways. First, every simulation
result carries its research boundary in its own identity material: the trace
binds a fixed Ring-4 boundary object, so the artifact that would be quoted as
evidence physically contains the statement of what it is not evidence for.
Second, the interpretation rule for a green fixture is stated normatively and
verbatim in the architecture's own documents and is repeated in this
monograph wherever a result is reported: the software performed the declared
transition; nothing else. Third, the evidence taxonomy (developed later in
the monograph) assigns software receipts, experiments, measurement
validation, safety evidence, and authority to disjoint categories, so that a
receipt cannot be counted in a category it does not belong to without a
detectable category error. The separation costs something --- it forbids the
architecture from ever advertising its most exciting-sounding
interpretations --- and that cost is accepted deliberately.

\subsection{Separation II: design value versus measurement}\label{ch01:sep:designmeasure}

\begin{ruleenv}[Design value versus measurement]\label{ch01:rule:designmeasure}
A locally supplied parameter or calibration label remains structural metadata
until authenticated measurement evidence and uncertainty are supplied. The
architecture never promotes a value from the design ledger to the measurement
ledger on the strength of anything the model author alone can construct.
\end{ruleenv}

The failure this separation prevents is \emph{self-certification}: the
laundering of an assumption into a datum by passing it through
official-looking machinery. Again an extended example makes the mechanism
concrete. A modeler declares a test point on a net and attaches a
calibration reference: an instrument identity, a protocol name, an
uncertainty figure, a measurement kind. All of these fields are well-formed;
the design-rule checker verifies their shape, their unit compatibility with
the net, and their presence where safety-marked paths require probes. A
simulation run then emits test-point observations, each carrying measurement
identifiers with parent/child derivations and the calibration reference.

Every artifact in this chain looks like measurement. There are instrument
fields, protocol fields, uncertainty figures, measurement IDs, and derivation
graphs. And yet nothing was measured: every field was typed in by the same
author who typed in the model, and the ``observation'' is a value computed by
the author's own declared transfer rules. If the architecture allowed these
records to flow into its evidence stream, it would have built an automated
assumption-laundering machine --- more dangerous than an informal diagram,
because the informal diagram at least looks informal. The hardware world
knows this failure intimately: a datasheet's typical values are design
targets until a characterization lab, with traceable instruments and stated
uncertainty, turns them into measurements, and no amount of typesetting the
target in the measurement column makes it one.

A second form of the same failure operates on parameters rather than
observations. A transfer-stage gain of $0.85$, a threshold of $0.9$, a
clamp of $0.1$ --- every such number in a board is a \emph{design value}: a
quantity the modeler chose, perhaps thoughtfully, perhaps arbitrarily, and
wrote into the model text. Design values are indispensable; nothing can be
checked or simulated without them. The failure is the citation loop: the
model is simulated with the chosen parameters, the simulation exhibits some
behavior, the behavior is written up, and in the write-up the parameters
have become findings --- ``the refusal stage saturates at $0.1$'' reported
in the tone of a discovered property rather than a typed-in choice. Six
months later the write-up is a source, and the design value has completed
its unearned journey to datum. The architecture's countermeasure is
ledger-keeping at the type level: parameters are typed as parameters
everywhere --- in the board, in the DRC report, in the trace --- and no
representation exists in which a parameter's provenance as a supplied value
is absent. The number can be quoted, but it cannot be quoted \emph{as a
measurement} without the quotation being checkably false against the
artifact it cites.

The architecture therefore enforces the separation structurally. Locally
constructed calibration records are typed as structural, test-only metadata:
they satisfy the checker's shape requirements, they enable simulation
fixtures, and they are permanently ineligible to confer authenticated
evidence status. Nothing the modeler can write in the model text --- no
field, no flag, no self-hash --- crosses the ledger boundary. Crossing it
requires an authenticated evidence pathway that the model author alone
cannot fabricate, and at the bound state of the work described in this
monograph, that pathway has not been exercised at all: there is no live
calibrated model, no instrument chain, and no measurement claim anywhere in
the admitted evidence. The measurement ledger is not merely separate; it is,
at present, deliberately empty.

\subsection{Separation III: analogy versus identity}\label{ch01:sep:analogy}

\begin{ruleenv}[Analogy versus identity]\label{ch01:rule:analogy}
Names borrowed from circuits do not make CAS quantities SI electrical
quantities. Physical voltage, current, resistance, charge, flux, and power
are outside the candidate unless a separately typed, instrumented physical
profile is proposed and validated. No such profile exists at the state of
work described in this monograph.
\end{ruleenv}

The failure this separation prevents is \emph{nominal capture}: the drift by
which a borrowed name gradually acquires the semantics of its origin. The
mechanism is linguistic before it is technical. A model component is called a
\emph{fuse} because its declared behavior --- open the path when a threshold
is crossed --- is shaped like a fuse's behavior. The name is apt, memorable,
and dangerous. Within a few pages of informal writing, the fuse ``blows,''
``protects the downstream stages,'' and ``guarantees'' something; within a
few documents, the model containing it is a ``safety architecture''; and at
the end of the slope someone is asserting, on the strength of a noun, that a
cognitive system has an overcurrent-protection property that no one has
defined, let alone measured. The same slope exists for every borrowed term:
gain, impedance, saturation, rail, ground. Each is a compressed description
of declared software behavior, and each invites its physical semantics to
follow the name through the door.

The engineering baseline that this monograph audits provides a vivid
demonstration of why the separation must be mechanical rather than
rhetorical. The baseline's published Cognitive Safety Board --- its flagship
safety-flavored design --- contains a red-stage saturation clamp of $0.1$
upstream of a fuse threshold of $0.9$. Under the declared clamp semantics
the fuse input can never exceed $0.1$, and the protective crossing is
unreachable: the fuse, as published, cannot trip. (The finding is stated
precisely, with page anchors and a formalized unreachability argument, later
in the monograph; Chapter~\ref{ch03:source} locates it in the source basis.)
A reader relying on the name \emph{fuse} and the label \emph{safety} would
have attributed a protective property to a design whose own parameters
forbid it. The design-rule checker, which propagates conservative bounds
through the declared path and flags unreachable safety thresholds, catches
the contradiction mechanically --- because it reasons about declared
semantics and never about names.

The separation also governs the quantitative vocabulary. CAS quantities in
this architecture carry their own units and signal kinds, checked under a
closed admissibility table; none of them is an SI electrical quantity, and
the checker treats an attempt to wire incompatible kinds together as an
error rather than a metaphor. If future work ever proposes that some CAS
quantity has a physical correlate, the architecture's position is that this
is a new, separately typed, instrumented physical profile requiring its own
validation program --- not a renaming of an existing model quantity. Until
such a profile exists and is validated, every electrical-sounding word in
this monograph is a typed pun, and the type system knows it.

\subsection{The separations jointly}\label{ch01:sep:jointly}

The three separations are stated independently but function as a system.
Model-versus-world polices the \emph{output} direction: what a computed
result may be said to show. Design-value-versus-measurement polices the
\emph{input} direction: what a supplied number may be said to be.
Analogy-versus-identity polices the \emph{vocabulary}: what a name may be
said to mean. A modeling practice that enforced only two of the three would
leak through the third: with the first two but not the third, names would
smuggle physics; with the first and third but not the second,
self-certified ``measurements'' would seed honest simulations with
counterfeit data; with the second and third but not the first, honest inputs
and honest names would still be spent on world-claims the software cannot
support. The monograph will return to this triangulation repeatedly, because
almost every design decision in the architecture --- from tagged exact
values to trace-bound research boundaries to the adapter's refusal to coerce
--- is most easily explained as the enforcement of one of these separations
at a specific interface.

\section{The six contributions}\label{ch01:contributions}

Against this background, the monograph develops six contributions. They are
the contributions of the current candidate --- a software research artifact
at a bound source state --- and each is claimed at the software-contract
level only.

\begin{enumerate}
\item \textbf{A typed board model} separating semantic state, geometry,
signal kind, unit, component contract, calibration metadata, and evidence
provenance. The model distinguishes state literals from ND-cells, keeps the
strictly positive simplex address of a mixed profile disjoint from both, uses
tagged exact type-plus-value identity for finite carrier values, and runs
all decimal arithmetic under one immutable, versioned context. Board
identity includes the complete canonical structure, so any structural edit
changes the board's digest.

\item \textbf{A deterministic design-rule checker (DRC) with fail-closed
diagnostics.} The DRC computes one deterministic report under a closed
35-code \texttt{PCBA8-*} catalog spanning model, semantics, parameters,
states, bias, nets, test points, calibration, feedback, safety, and
TKS-adapter classes. Errors fail the board; warnings never become an
implicit clean result; unsupported or resource-bounded analysis yields an
unresolved diagnostic, never a pass.

\item \textbf{A narrow trace-emitting reference simulator} for declared
deterministic component profiles. Supported components declare ports,
transfer rules, clipping, statefulness, and time semantics; every run emits
a canonical trace binding board, component, DRC-report, catalog,
decimal-context, profile, and source identities together with typed samples,
events, observations, measurement derivations, and the fixed research
boundary.

\item \textbf{A source- and profile-bound adapter} from certified finite
component tables to the TKS v8 finite reference kernel. Only a certified
total, finite, deterministic, stateless, non-stochastic transition may cross;
the adapter returns an immutable certificate and proves at most type-correct
representation, carrying no automatic semantic theorem.

\item \textbf{A source audit that preserves, rather than silently repairs,
adverse findings.} The audit records that the baseline's consolidated
grammar does not cover its own published five-board corpus (four direct
syntactic counterexamples) and that the published Cognitive Safety Board's
protective crossing is unreachable under its declared parameters; both
findings are retained with page anchors as permanent fixtures.

\item \textbf{An evidence discipline} distinguishing green software tests
from experiments, measurement validation, safety evidence, and authority,
with every claim record at the bound state classified Ring~4 and
\textsc{not licensed}.
\end{enumerate}

Two remarks about this list. First, none of the six is an empirical result
about cognition; the list is deliberately and entirely architectural.
Second, the fifth contribution deserves emphasis because it is the one most
often absent from projects of this kind: the architecture treats the
discovery that its own source material is defective --- syntactically
incomplete in its grammar, self-contradictory in its flagship safety design
--- as a result to be formalized and guarded, not an embarrassment to be
patched. A checker that has flagged its own canon is a checker whose
fail-closed behavior has been demonstrated on a case that mattered.

\section{The thesis: a qualified affirmative at the software-contract level}\label{ch01:thesis}

\begin{hypothesis}[Thesis of the monograph]\label{ch01:hyp:thesis}
The research question of Section~\ref{ch01:question} admits a qualified
affirmative answer at the software-contract level: a circuit-shaped modeling
language, equipped with the type system, design-rule checker, trace-bound
simulator, adapter boundaries, source audit, and evidence discipline
developed in this monograph, makes cognitive-system hypotheses more explicit,
type-safe, replayable, and falsifiable than their informal-diagram
equivalents, while mechanically withholding physical, causal, safety, and
product claims.
\end{hypothesis}

Every word of the qualification is intended. \emph{Qualified}: the
affirmative is supported by bounded software evidence --- at the bound
source state, an integrated receipt records three green focused suites
totaling 208 test methods and 272 parameterized subtests, of which 82
methods belong to CAS PCBA v8 itself --- and by two formalized adverse
findings against the source; it is not supported by any experiment,
measurement, replication, or external review, all of which remain open.
\emph{At the software-contract level}: the affirmative asserts that the
software honors its declared contracts on its tested surface, and asserts
nothing about any system outside the software. The thesis would be
falsified by, for example, a demonstration that the checker passes a board
violating a cataloged rule, that the simulator's traces are not
deterministic under a declared profile, that a coercion crosses a boundary
the type system claims to close, or that the preserved adverse findings
misstate the source at the cited anchors. The thesis is \emph{not}
strengthened by any accumulation of green fixtures beyond its tested
surface, and the monograph will not argue as if it were.

\subsection{Falsifiability surfaces of the thesis}\label{ch01:thesis:falsifiers}

Because the monograph asks its models to expose failure surfaces, it owes
the same of its thesis. The qualified affirmative of
Hypothesis~\ref{ch01:hyp:thesis} is falsifiable on at least five distinct
surfaces, each checkable by a party holding the bound artifacts, and the
monograph invites exactly this scrutiny:

\begin{enumerate}
\item \textbf{Checker soundness.} Exhibit a board that violates a rule of
the closed diagnostic catalog and nonetheless receives a clean report
from the checker at the bound implementation. One such witness defeats
the fail-closed claim for the violated rule class.
\item \textbf{Determinism.} Exhibit two runs of the checker or the
simulator, on identical inputs under identical declared profiles and
contexts, producing non-identical reports or traces. One such pair
defeats the replayability claim.
\item \textbf{Boundary leakage.} Exhibit a coercion the type system
claims to forbid --- a state literal accepted where a simplex address is
required, an untagged value crossing the adapter, a locally constructed
calibration record acquiring authenticated evidence status --- actually
succeeding in the bound implementation.
\item \textbf{Audit fidelity.} Show that a preserved adverse finding
misstates the source at its cited anchors --- that the consolidated
grammar in fact contains an alternative producing a cited surface form,
or that the safety-board listing does not carry the cited clamp and
threshold values at the cited pages of the hashed build.
\item \textbf{Evidence typing.} Show that the architecture's own
artifacts, read as typed, license an inference the claim quarantine
prohibits --- that some receipt, certificate, or report, under its
declared semantics rather than a careless reading, constitutes a
physical, causal, safety, or product claim.
\end{enumerate}

Symmetrically, the thesis is \emph{not} falsified by demonstrations that
belong to the withheld levels: showing that no board in the candidate
corresponds to any real cognitive mechanism would defeat a validity
claim the thesis does not make, and showing that the checker's rules
encode contestable modeling taste would defeat an objectivity claim it
does not make either. The thesis lives and dies at the software
contract, on the five surfaces above.

It is worth stating plainly what would be required to answer the research
question affirmatively at any level beyond the software contract, because
the distance is part of the honest content of this work: a complete
versioned grammar with round-trip conformance; stateful and stochastic
component semantics; live calibration with authenticated instruments and
uncertainty; preregistered held-out validation; independent replication in
a hermetic toolchain; and separate formal, systems, experimental, safety,
and claim-language review. That program is laid out, dependency-ordered, at
the end of the monograph. None of it has begun in the admitted evidence,
and nothing in the intervening chapters will pretend otherwise.

\section{What is gained over informal practice}\label{ch01:gained}

Because the thesis withholds so much, it is fair to ask what, concretely,
is left --- what a modeler gains by writing a cognitive-system hypothesis
as a checked board rather than as a diagram. The answer, developed
mechanism by mechanism in the later chapters, is that whole categories of
error become \emph{observable earlier}, mechanically, at design time
rather than at argument time. Eight categories recur throughout the
monograph and are worth previewing as a list, since together they are the
practical content of the qualified affirmative:

\begin{enumerate}
\item \textbf{Scale and unit substitution.} Quantities of different
declared kinds or units cannot be wired together; the admissibility table
is closed, and a mismatch is a distinct diagnostic rather than a silent
reinterpretation.
\item \textbf{Ambiguous semantic crosswalks.} Semantic references live in
disjoint namespaces with source and definition hashes; two objects sharing
a label or a position never silently become the same object, so an
equivocation must be written down as an explicit (and checkable)
translation to exist at all.
\item \textbf{Incomplete grammar claims.} A language that claims to
describe a corpus can be audited, construct by construct, against that
corpus --- and when the monograph's own source canon failed such an audit,
the failure was recordable at page-anchor resolution precisely because the
claim had been made mechanical.
\item \textbf{Multiple-driver and unresolved-port defects.} The netlist
discipline makes contention, floating inputs, and dangling references
findings rather than latent ambiguities.
\item \textbf{Unreachable threshold paths.} Conservative bound propagation
can prove that a declared protective crossing cannot occur --- the
category of defect that condemns the source's own safety board.
\item \textbf{Undisciplined feedback.} Cycles must be declared, delayed,
calibrated, and bounded; an undeclared cycle is an error, and an
unanalyzable one is unresolved, never clean.
\item \textbf{Fabricated or disconnected measurement lineage.} Measurement
records must form acyclic, existence-checked derivation chains bound to
calibration references, and locally asserted calibration can never
self-authenticate.
\item \textbf{Receipts misrepresented as authority.} Software receipts,
claim records, and review ledgers are typed artifacts; a green suite, a
self-hash, or a generated receipt is checkably \emph{not} a license, a
certification, or a review.
\end{enumerate}

None of these categories requires the cognitive hypotheses to be true.
That is the design: the gains are gains in the hygiene of hypothesizing,
available in full even if every hypothesis so expressed turns out false.
An informal diagram can be wrong invisibly in all eight ways at once; a
checked board can be wrong too --- about the world, richly and perhaps
irreparably --- but not invisibly in these ways, and the difference is
what the word ``falsifiable'' in the research question is doing.

\section{Claim discipline and artifact status}\label{ch01:discipline}

\begin{boundary}[Standing claim boundary for the monograph]\label{ch01:bnd:standing}
The work described here is a research candidate with artifact status
\textsc{draft / open / ring 4 / not licensed / held}. Ring~4 designates the
outermost research ring: no live, scientific, canon, product, license, or
disclosure authority exists or is created by this text. All claim records at
the bound source state are Ring~4 and \textsc{not licensed}. Drafting,
compiling, or circulating this monograph promotes no claim; independent
systems, safety, regulatory, clinical, and formal review remains open.
\end{boundary}

The following claim quarantine, inherited verbatim in substance from the
governing research paper and the Gate~4 audit, applies to every chapter. The
items below remain hypotheses or prohibited product inferences, and no
sentence in this monograph asserts any of them: exact thought/transistor
isomorphism; shared four-region identity; overload-induced semantic
involution; Fisher--Rao causal gain; physical electrical identity for CAS
quantities; semantic correctness inferred from a response-region label;
safety effectiveness, certification, regulatory compliance, or clinical use;
complete compiler conformance; production readiness; and disclosure or
license authority derived from a self-hash or generated receipt. Historical
AI-authored, self-hashed licensing receipts are not an authenticated
reviewer or legal authority, and none is treated as one here.

Three further conventions govern the text. First, \emph{background is
background}: Chapter~\ref{ch02:background} surveys established engineering
and mathematical literature, and nothing in that survey is presented as a
novelty of this project or as validation of CAS. Second, \emph{sources are
bound by identity}: every fact about the engineering baseline and the
controlling repository artifacts is cited against the exact content hashes
recorded in Chapter~\ref{ch03:source}, and where the sources do not specify
a fact, the monograph says so rather than estimating. Third, \emph{negative
results are first-class}: adverse findings are stated in their own
environments, at full evidential resolution, and are never softened in
summary.

\section{Reading guide}\label{ch01:guide}

The monograph is organized so that a reader can locate the boundary of every
claim without reading linearly, but a linear reading is the intended one.

Chapter~\ref{ch02:background} is a tutorial survey of the background this
work stands on: hardware description languages and what VHDL demonstrated
about formal, machine-readable design description; Statecharts and the
explicit semantics of state, hierarchy, and concurrency; the industrial
practice of netlists, design-rule checking, and electrical-rule-style
connectivity checks; graph algorithms for circuit analysis; the
verification/validation distinction; property-based testing; provenance
modeling; floating-point versus decimal arithmetic control; and a primer on
information geometry over probability simplices, including the Fisher--Rao
metric and its closed-form geodesic distance. The chapter is deliberately
self-contained so that readers from cognitive-science, software, and
hardware backgrounds can each acquire the vocabulary the later chapters
assume. Nothing in it is new, and it is careful to say so.

Chapter~\ref{ch03:source} fixes the versioned source basis: the 743-page
engineering baseline with its source ID and SHA-256 hash; the reasons the
work identifies sources by content hash rather than by an invented public
version; the baseline's own disclaimer and the Phase-0 authority matrix;
the controlling repository artifacts and their hashes; the downstream
read-only boundaries into the TKS v8 finite kernel and the Noetic Topos 2.0
evidence discipline; and the read-only status of the Gate-5 preregistration
kernel. Everything later in the monograph that mentions a source fact
resolves through this chapter.

The chapters that follow the source basis develop the architecture in
dependency order: the typed semantic and board model, including state
types, tagged values, the single decimal context, and board identity; the
design-rule checker and its closed diagnostic catalog, check class by check
class; the narrow deterministic simulator and its canonical trace bindings;
the geometry profiles --- the Fisher--Rao baseline and the optional
projector profile --- with their disclosed numerical limits; and the
certified-transition adapter into the TKS v8 finite kernel, together with
the Noetic Topos ambient/resolved boundary. After the architecture, the
monograph presents the adverse source audit at full resolution --- the
grammar-coverage finding with its board-by-board evidence table, and the
Cognitive Safety Board finding with the formalized unreachability argument
--- followed by the bounded software evidence at the bound commit, the
threats to validity, and the dependency-ordered research program. A reader
interested only in what this work does \emph{not} claim can read
Section~\ref{ch01:discipline}, the boundary environments of each chapter,
and the threats-to-validity chapter, and will have an accurate picture.

A note on tone. The monograph is written in the declarative voice of an
engineering document, and long stretches of it describe machinery that
works as declared. The reader is asked to hold that voice together with the
standing boundary of Section~\ref{ch01:discipline} at all times: ``works''
always means ``performs its declared software contract on its tested
surface,'' the tested surface is finite and disclosed, and the most
important results in the monograph are the two places where the machinery,
applied to its own source canon, correctly said \emph{no}.
